\chapter{Visibilities and Storage}
\label{chap:visibilities_and_storage}

Now that we have timing and positioning figured out, it's time to compute visibility data.
We re-pfb our channelized data before applying a cross-correlation, explained in XXXX.
We save data as a pyuv format, which is detailed in XXXX.

\section{Re-PFB}

Main idea is to take our baseband data. It is already channelized, but coarsely. 
We Inverse PFB it, and return to (approximately) the voltage-time data.
We then apply another PFB, with some oversampling constant.
Get finer channels, and also specify a certain region we want to look at.
We then cross-correlate and get our visibilities with the oversampled channel resolution.
This is a better method because....

In our case, we use these numbers XXXX


\section{PYUV}

%intro sentence
The cross correlation gives us plenty of data.
%motivation
We save everything in the specific PYUV format which allows for saving several baselines and spectral windows,
while having relatively simple handling and updating.

Below is an outline of the main objects in such files, and the conventions we use for them.


\subsection{Constants}


\begin{table}[H]
\centering
\begin{tabular}{|l|p{15cm}|}
\hline
\textbf{Field} & \textbf{Description} \\
\hline
\texttt{Nbls} & Number of baselines. \\
\hline
\texttt{Ntimes} & Number of times. \\
\hline
\texttt{Nblts} & Number of baseline-times. Not necessarily equal to \texttt{Nbls} $\times$ \texttt{Ntimes}. \\
\hline
\texttt{Nfreqs} & Number of frequency channels. \\
\hline
\texttt{Npols} & Number of polarizations. \\
\hline
\texttt{Nspws} & Number of spectral windows. \\
\hline
\texttt{Nphase} & Number of phase centers in the data. \\
\hline
\texttt{Nants\_data} & Number of antennas with data present (i.e., number of unique entries in \texttt{ant\_1\_array} and \texttt{ant\_2\_array}). May be smaller than the number of antennas in the array. \\
\hline
\texttt{vis\_units} & Visibility units. Options are: \texttt{"uncalib"}, \texttt{"Jy"}, or \texttt{"K str"}. \\
\hline
\end{tabular}
\caption{Key PYUV parameters. Will determine much of the rest.}
\end{table}


\subsection{Data}

\subsubsection{Visibility-Related}

Explain how baseline-times are arranged??? How to extract data??

\begin{table}[H]
\centering
\begin{tabular}{|l|l|l|p{5cm}|}
\hline
\textbf{Name} & \textbf{Shape} & \textbf{Type} & \textbf{Description} \\
\hline
\texttt{data\_array} & (Nblts, Nfreqs, Npols) & complex float & Visibility data in units of \texttt{vis\_units} \\
\hline
\texttt{flag\_array} & (Nblts, Nfreqs, Npols) & bool & True indicates flagged data \\
\hline
\texttt{nsample\_array} & (Nblts, Nfreqs, Npols) & float & Number of samples included in the visibility chunk\footnotemark[1] \\
\hline
\end{tabular}
\caption{Where the visibility information lives.}
\end{table}

\footnotetext[1]{Also often used to check for flagging.}

\subsubsection{Time-Related}
\begin{table}[H]
\centering
\begin{tabular}{|l|l|l|l|}
\hline
\textbf{Name} & \textbf{Shape} & \textbf{Units} & \textbf{Description} \\
\hline
\texttt{time\_array} & (Nblts) & Julian Date & Time at center of integration. \\
\hline
\texttt{lst\_array} & (Nblts) & radians & Local Apparent Sidereal Time (LAST). \\
\hline
\texttt{integration\_time} & (Nblts) & seconds & Duration of each visibility integration. \\
\hline
\end{tabular}
\caption{Time-related data. Required to reconstruct time information from \texttt{data\_array}}
\end{table}

\subsubsection{Baseline-Related}
\begin{table}[H]
\centering
\begin{tabular}{|l|l|l|l|}
\hline
\textbf{Name} & \textbf{Shape} & \textbf{Type} & \textbf{Description} \\
\hline
\texttt{ant\_1\_array} & (Nblts) & int & First antenna in baseline. \\
\hline
\texttt{ant\_2\_array} & (Nblts) & int & Second antenna in baseline. \\
\hline
\texttt{baseline\_array} & (Nblts) & int & Encodes ant1 and ant2 (we use: $2048 \times \texttt{ant1} + \texttt{ant2} + 2^{16}$)\footnotemark[2]. \\
\hline
\end{tabular}
\caption{Antenna and baseline data. Also required to reconstruct baseline information from \texttt{data\_array}}
\end{table}

\footnotetext[2]{This special code is for compatibility with other data formats, and for efficiency in sorting methods.}

\subsubsection{Frequency and Polarization}
\begin{table}[H]
\centering
\begin{tabular}{|l|l|l|l|}
\hline
\textbf{Name} & \textbf{Shape} & \textbf{Units/Type} & \textbf{Description} \\
\hline
\texttt{freq\_array} & (Nfreqs) & Hz & Center frequency of each channel. \\
\texttt{channel\_width} & (Nfreqs) & Hz & Width of each frequency channel. \\
\texttt{polarization\_array} & (Npols) & int & Polarization codes.\footnotemark[3] \\
\hline
\end{tabular}
\caption{Spectral and polarization data.}
\end{table}

\footnotetext[3]{Different integers serve to represent which exact polarization we are working with. 
Will not go through it here, see AIPS Memo 117}


\subsubsection{Spectral Window}
\begin{table}[H]
\centering
\begin{tabular}{|l|l|l|l|}
\hline
\textbf{Name} & \textbf{Shape} & \textbf{Type} & \textbf{Description} \\
\hline
\texttt{spw\_array} & (Nspws) & int & Spectral window identifiers. \\
\texttt{flex\_spw\_id\_array} & (Nfreqs) & int & Maps frequencies to spectral windows. \\
\hline
\end{tabular}
\caption{Spectral window information.}
\end{table}



\section{Storage System}

How the pyuv files work (includingconventions for storage), 
what the storage system is, and how to access them (with examples!).

